\chapter{组合数学}
\label{chap:combinatorics}

第~\ref{part:probability-stochastic-processes-statistics}部分用分布描述不确定结果，
但许多学习问题还要求从有限候选中作出选择。
标签是否启用、规则是否满足、特征是否保留，最终都可写成有限个离散变量的赋值。
候选有限只保证可以逐个列出，并不说明候选数量可以承受，也不说明每个候选都符合规则，
更不说明最优候选能在可接受的时间内找到。要区分这些问题，需要分别研究计数、约束、
目标结构和计算资源。

本章沿用同一组二值变量。令
\[
  V=\{1,\ldots,n\},\qquad
  \bm{z}=(z_i)_{i\in V}\in\{0,1\}^n,
  \qquad
  S(\bm{z})=\{i\in V:z_i=1\}.
\]
$z_i=1$表示选择第$i$个候选，集合$S(\bm{z})$记录完整选择。贯穿各节的小例子取
$V_0=\{1,2,3,4\}$：候选$1$与$2$这两种基础标签必须恰选一个，候选$3$依赖候选$1$，
候选$4$则要求候选$2$或$3$至少出现一个。计数回答可能有多少个$S$，逻辑规则从中
筛出可行集合，集合函数比较这些可行集合，而计算理论追问相应程序能否停机以及资源怎样
随输入增长。变量之间的关系怎样支持局部计算，将由第~\ref{chap:graph-theory}章继续讨论。

\section{集合族与组合计数}
\label{sec:comb-families-counting}

\subsection{子集族把允许的候选组织起来}

第~\ref{chap:sets-functions-and-proofs}章已经给出幂集与组合数。这里关心的不是孤立的
计数公式，而是限制改变后候选规模怎样变化。一个\term{集合族}（set family）
$\mathcal{F}\subseteq\mathcal{P}(V)$是一组允许讨论的子集。若不加任何限制，
$\mathcal{F}=\mathcal{P}(V)$；每个元素独立地选择或不选择，所以
\begin{equation}
  |\mathcal{P}(V)|=2^n.
  \label{eq:comb-power-set-count}
\end{equation}
这个结论也可由二项式定理得到：
$\sum_{k=0}^{n}\binom{n}{k}=(1+1)^n=2^n$。两种证明分别强调二值赋值与按集合大小
分层，它们计算的是同一批对象。

同样是从$n$个不同元素中取出$k$个，是否记录顺序会改变计数对象。若依次填入
$k$个位置，其中$0\leq k\leq n$，每个位置不能重复使用元素，方案数为
\[
  n(n-1)\cdots(n-k+1)=\frac{n!}{(n-k)!}.
\]
若只记录被选中的子集，每个$k$元子集会被它的$k!$种排列重复表示，因而方案数为
$\binom{n}{k}=n!/[k!(n-k)!]$。例如从$\{1,2,3,4\}$中取两个元素，有$12$个
有序对，却只有$6$个二元子集；$(1,3)$与$(3,1)$在前一个问题中不同，在后一个问题中
表示同一个选择。计数公式必须随“什么结果算作同一个”一起给出。

每个$S\subseteq V$也确定一个有序二分划分$(S,V\setminus S)$，所以有序二分划分
仍有$2^n$个。若两部分没有名称，$(S,V\setminus S)$与$(V\setminus S,S)$表示
同一个划分；当$n\geq1$时没有子集等于自己的补集，每个无序划分恰被计算两次，数量为
$2^{n-1}$。是否区分两侧决定了计数对象，不能只看它们都把$V$分成两部分。
这里允许某一部分为空；若要求两部分都非空，有序与无序数量应分别减为
$2^n-2$与$2^{n-1}-1$。当$n=0$时，交换两侧不会产生两个不同表示，上述除以$2$的
论证也不适用。

若只允许恰好选择$k$项，集合族变为
\[
  \binom{V}{k}=\{S\subseteq V:|S|=k\},
  \qquad
  \left|\binom{V}{k}\right|=\binom{n}{k}.
\]
若允许至多选择$k$项，则数量为$\sum_{j=0}^{k}\binom{n}{j}$。因此“候选总数有限”
仍可能掩盖完全不同的增长规模：固定$k$时，该和关于$n$至多按$n^k$增长；
$k$随$n$线性增长时，候选数却可能接近$2^n$。

计数常可通过一个变量把问题拆成两个较小问题。记$a_{n,k}$为从$n$个元素中选择
$k$个的方案数。对$n\geq1$，固定元素$n$后，所有方案被不重不漏地分成“不含$n$”
与“含$n$”两类。为使边界项也有定义，约定$a_{0,0}=1$，并在$k<0$或$k>n$时令
$a_{n,k}=0$。于是对$n\geq1$和任意整数$k$都有
\begin{equation}
  a_{n,k}=a_{n-1,k}+a_{n-1,k-1}.
  \label{eq:comb-binomial-recurrence}
\end{equation}
这既是Pascal恒等式，也是一个递推算法。若只保存上一行，计算全部
$a_{n,0},\ldots,a_{n,n}$需要$O(n^2)$次加法和$O(n)$个存储单元。
闭式公式、递推关系和计算过程由此承担不同作用：闭式便于分析规模，递推揭示分组结构，
算法还必须说明求值顺序和资源。

\subsection{二值标记的组合计数}

第~\ref{chap:functional-analysis-and-approximation}章式~
\eqref{eq:functional-growth-function}已经定义增长函数，并用有限输入上的限制来刻画
二值函数类的组合丰富度。这里把每种标记改写成一个子集，以便直接使用本节的计数方法。
给定输入$x_1,\ldots,x_m$和二值函数类$\mathcal{H}$，每个$h\in\mathcal{H}$确定
\[
  A_h=\{i\in\{1,\ldots,m\}:h(x_i)=1\}.
\]
函数在这组输入上的标记向量与$A_h$一一对应，因此不同标记的数量就是集合族
$\{A_h:h\in\mathcal{H}\}\subseteq\mathcal{P}(\{1,\ldots,m\})$的大小。式~
\eqref{eq:comb-power-set-count}随即给出上界$2^m$；增长函数只是在所有输入组中取这个
数量的上确界。

这个表示还能直接计算受基数约束的函数族。令$\mathcal{X}$为无限集，固定$k\geq0$，并设
\[
  \mathcal{H}_k
  =
  \{h_S:S\subseteq\mathcal{X},\ |S|\leq k\},
  \qquad
  h_S(x)=\mathbb{I}\{x\in S\}.
\]
对$m$个互异输入，$\mathcal{H}_k$能够实现的正类索引集合恰是所有大小不超过$k$的
子集，因而
\[
  \Pi_{\mathcal{H}_k}(m)
  =
  \sum_{j=0}^{\min\{k,m\}}\binom{m}{j}.
\]
固定$k$时，这个数量关于$m$至多按$m^k$增长；当$k\geq m$时，它恢复为$2^m$。
同一个幂集计数由此区分出多项式规模与指数规模。它仍只描述有限输入上的模式数，
不直接给出统计误差界，也不保证能够高效找到实现某种标记的函数。

一个可逐项复核的例子是实数轴上的区间分类器。令
$x_1<\cdots<x_m$互异，函数$h_{a,b}(x)=\mathbb{I}\{a\leq x\leq b\}$把落在某个
闭区间内的输入标为$1$。非空正类索引必为连续块
$\{i,i+1,\ldots,j\}$，这样的块由$1\leq i\leq j\leq m$唯一确定；再加上空正类，
可实现模式数恰为
\[
  1+\sum_{i=1}^{m}(m-i+1)
  =
  1+\frac{m(m+1)}2.
\]
当$m=3$时，七个模式为$000,100,010,001,110,011,111$，只有$101$无法实现。
因此三个有序点不能被该函数类打散；两个点的四种模式都能实现。这个算例也说明，
参数$a,b$取值无限多，并不意味着有限输入上的标记模式也有无限多个。

\subsection{容斥原理校正重复计数}

直接把若干坏情形的数量相加会重复计算同时违反多条规则的赋值。
\term{容斥原理}（inclusion--exclusion principle）用交集项逐层校正这种重复。

\begin{theorem}[有限容斥原理]
\label{thm:comb-inclusion-exclusion}
设$A_1,\ldots,A_m$是有限集合。则
\begin{equation}
  \left|\bigcup_{i=1}^{m}A_i\right|
  =
  \sum_{\varnothing\neq I\subseteq\{1,\ldots,m\}}
  (-1)^{|I|+1}
  \left|\bigcap_{i\in I}A_i\right|.
  \label{eq:comb-inclusion-exclusion}
\end{equation}
\end{theorem}

\begin{proof}
固定任意元素$x$，设它恰好属于$r$个集合。若$r=0$，它在等式两侧都没有贡献。
若$r\geq1$，右侧从含有$x$的$r$个集合中选择非空索引集，给$x$计算的总系数为
\[
  \sum_{k=1}^{r}(-1)^{k+1}\binom{r}{k}
  =
  1-\sum_{k=0}^{r}(-1)^k\binom{r}{k}
  =1-(1-1)^r=1.
\]
所以并集中的每个元素在右侧恰好被计算一次，对所有$x$求和即得结论。
\end{proof}

回到$V_0$上的四个变量。把三条规则的违反事件记为
\[
  \begin{aligned}
    A_1&=\{\bm{z}:z_1=z_2\},\\
    A_2&=\{\bm{z}:z_3=1,\ z_1=0\},\\
    A_3&=\{\bm{z}:z_4=1,\ z_2=0,\ z_3=0\}.
  \end{aligned}
\]
它们分别表示基础标签没有恰选一个、缺少候选$1$时仍选择候选$3$，以及候选$4$缺少前提。
全部赋值有$16$个，并且
\[
  |A_1|=8,\quad |A_2|=4,\quad |A_3|=2.
\]
交集大小为
\[
  |A_1\cap A_2|=2,\qquad
  |A_1\cap A_3|=1,\qquad
  |A_2\cap A_3|=|A_1\cap A_2\cap A_3|=0.
\]
由式~\eqref{eq:comb-inclusion-exclusion}，坏赋值共有
$8+4+2-2-1=11$个，因此可行赋值有$16-11=5$个。它们对应
\begin{equation}
  \mathcal{C}_0
  =
  \bigl\{
  \{1\},\{1,3\},\{1,3,4\},\{2\},\{2,4\}
  \bigr\}.
  \label{eq:comb-running-feasible-family}
\end{equation}
这个计算在第~\ref{sec:comb-logic-constraints}节还会由逻辑公式复核。

\subsection{子集格上的 Möbius 反演}

容斥处理的是若干集合的重叠。一个更一般的问题是：若已知每个$S$上的总量由其全部
子集贡献累加而成，能否还原每个子集自身新增的贡献？子集按包含关系形成偏序，
\term{Möbius反演}（Möbius inversion）给出这一累加变换的逆。

\begin{theorem}[子集格上的 Möbius 反演]
\label{thm:comb-mobius-inversion}
设$g,h:\mathcal{P}(V)\to\mathbb{R}$满足
\begin{equation}
  g(S)=\sum_{T\subseteq S}h(T)
  \qquad\text{对每个 }S\subseteq V.
  \label{eq:comb-zeta-transform}
\end{equation}
则
\begin{equation}
  h(S)
  =
  \sum_{T\subseteq S}(-1)^{|S|-|T|}g(T).
  \label{eq:comb-mobius-inversion}
\end{equation}
反过来，由式~\eqref{eq:comb-mobius-inversion}定义$h$也会恢复
式~\eqref{eq:comb-zeta-transform}。
\end{theorem}

更一般的有限偏序集把$T\subseteq S$换成$T\preceq S$。其Möbius函数由
$\mu(S,S)=1$以及
$\mu(T,S)=-\sum_{T\preceq R\prec S}\mu(T,R)$递归确定，反演系数不再只由
$|S|-|T|$决定。本章只使用子集格，此时递归解恰为
$\mu(T,S)=(-1)^{|S|-|T|}$。

\begin{proof}
把式~\eqref{eq:comb-zeta-transform}代入反演式右侧并交换有限求和次序，得到
\begin{align*}
  \sum_{T\subseteq S}(-1)^{|S|-|T|}g(T)
  &=
  \sum_{T\subseteq S}(-1)^{|S|-|T|}
  \sum_{R\subseteq T}h(R)\\
  &=
  \sum_{R\subseteq S}h(R)
  \sum_{T:R\subseteq T\subseteq S}
  (-1)^{|S|-|T|}.
\end{align*}
固定$R$后，中间集合$T$由$S\setminus R$中哪些元素被加入决定。若
$q=|S\setminus R|$，内层和等于
\[
  \sum_{j=0}^{q}\binom{q}{j}(-1)^{q-j}
  =(1-1)^q.
\]
它在$R=S$时为$1$，在$R\neq S$时为$0$，所以全部较小子集的贡献抵消，只留下
$h(S)$。同一抵消关系也证明反向代入恢复$g$。
\end{proof}

在小集合$V_0$上，假设累计得分具有形式
\[
  g(S)
  =
  2|S|
  +3\mathbb{I}\{\{1,3\}\subseteq S\}
  -\mathbb{I}\{\{2,4\}\subseteq S\}.
\]
每个单项贡献$2$，组合$\{1,3\}$额外贡献$3$，而组合$\{2,4\}$产生$-1$的交互。
四个元素共有$16$个子集。把每个子集的累计值与反演恢复的局部贡献并列，可以完整核对
正向累加与反向消去：
\begin{table}[htbp]
  \centering
  \caption{四元素子集格上的累计值与局部贡献。除四个单项、交互
  $\{1,3\}$和$\{2,4\}$外，其余局部贡献均为零。}
  \label{tab:comb-four-element-mobius}
  \begin{tabular}{ccc@{\qquad}ccc}
    \toprule
    $S$ & $g(S)$ & $h(S)$ & $S$ & $g(S)$ & $h(S)$ \\
    \midrule
    $\varnothing$ & $0$ & $0$ & $\{1,2\}$ & $4$ & $0$ \\
    $\{1\}$ & $2$ & $2$ & $\{1,3\}$ & $7$ & $3$ \\
    $\{2\}$ & $2$ & $2$ & $\{1,4\}$ & $4$ & $0$ \\
    $\{3\}$ & $2$ & $2$ & $\{2,3\}$ & $4$ & $0$ \\
    $\{4\}$ & $2$ & $2$ & $\{2,4\}$ & $3$ & $-1$ \\
    $\{3,4\}$ & $4$ & $0$ & $\{1,2,3\}$ & $9$ & $0$ \\
    $\{1,2,4\}$ & $5$ & $0$ & $\{1,3,4\}$ & $9$ & $0$ \\
    $\{2,3,4\}$ & $5$ & $0$ & $V_0$ & $10$ & $0$ \\
    \bottomrule
  \end{tabular}
\end{table}
若只观察各子集的累计值，反演给出
\[
  h(\{i\})=g(\{i\})-g(\varnothing)=2,
\]
以及
\[
  \begin{aligned}
  h(\{1,3\})
  &=g(\{1,3\})-g(\{1\})-g(\{3\})+g(\varnothing)=3,\\
  h(\{2,4\})
  &=g(\{2,4\})-g(\{2\})-g(\{4\})+g(\varnothing)=-1.
  \end{aligned}
\]
其余大小至少为$2$的局部贡献均为零。反演并没有假设贡献为正；它只保证累计量与局部
交互之间可以互相重构。

例如对$S=\{1,2,3\}$，累计值为$g(S)=9$，三个二元子集的累计值依次为
$4,7,4$，三个单元素累计值之和为$6$，且$g(\varnothing)=0$。反演确实给出
\[
  h(\{1,2,3\})=9-(4+7+4)+6-0=0.
\]
把已经恢复的单项贡献$2,2,2$与二元贡献$3$重新相加，也得到
$g(\{1,2,3\})=9$。正向累加与反向消去由此可以在同一组数值上相互复核。

容斥原理可以看成同一反演机制作用于指示函数。正负号
$(-1)^{|S|-|T|}$不是需要另行记忆的巧合，而是“所有严格更小子集相互抵消”所需的系数。
后续无向图模型把对数联合量写成各变量子集的局部交互时，会复用这个结论；图结构承担的
额外工作，是证明不构成团的子集交互为零。Möbius反演本身只是一条有限集合恒等式，
不会自动产生图分解，也不保证所得局部项非负。

\section{逻辑与约束推理}
\label{sec:comb-logic-constraints}

\subsection{公式把规则变成可检验条件}

计数把每个$\bm{z}\in\{0,1\}^n$都视为候选，实际任务却常只允许其中一部分。
\term{布尔公式}（Boolean formula）用命题变量和逻辑连接词描述允许的赋值。
对命题$P,Q$，合取$P\land Q$要求两者都真，析取$P\lor Q$要求至少一个为真，
否定$\neg P$翻转真假，蕴含$P\Rightarrow Q$则禁止$P$真而$Q$假的情形。
蕴含可以改写为
\begin{equation}
  P\Rightarrow Q
  \quad\Longleftrightarrow\quad
  \neg P\lor Q.
  \label{eq:comb-implication-clause}
\end{equation}

形式上，\term{赋值}（assignment）是映射
$\nu:\{z_1,\ldots,z_n\}\to\{\mathrm{false},\mathrm{true}\}$；连接词的真值规则
递归地确定$\nu$是否满足一个公式。满足公式$\Phi$的赋值也称为$\Phi$的
\term{模型}（model），全体模型记为$\operatorname{Mod}(\Phi)$。这个“模型”只指
逻辑公式的满足赋值，不等同于后续章节中由参数描述的统计模型。

否定复合命题时，De Morgan律给出
\begin{equation}
  \neg(P\land Q)\Longleftrightarrow(\neg P\lor\neg Q),
  \qquad
  \neg(P\lor Q)\Longleftrightarrow(\neg P\land\neg Q).
  \label{eq:comb-de-morgan}
\end{equation}
以第一式为例，$(P,Q)$依次取$(0,0),(0,1),(1,0),(1,1)$时，等式两侧的真值都依次为
$1,1,1,0$；第二式同理得到$1,0,0,0$。逐行相同证明两边具有相同模型集合，而不是
仅在某个选定赋值上碰巧相等。

把变量$z_i$同时看作命题“候选$i$被选择”，贯穿例子的规则可写成
\begin{equation}
  \Phi(\bm{z})
  =
  (z_1\lor z_2)
  \land(\neg z_1\lor\neg z_2)
  \land(\neg z_3\lor z_1)
  \land(\neg z_4\lor z_2\lor z_3).
  \label{eq:comb-running-cnf}
\end{equation}
前两个子句共同表达$z_1,z_2$恰有一个为真；第三个子句表达
$z_3\Rightarrow z_1$；第四个子句表达$z_4\Rightarrow(z_2\lor z_3)$。
像式~\eqref{eq:comb-running-cnf}这样由若干析取子句合取而成的公式称为
\term{合取范式}（conjunctive normal form, CNF）。

赋值$\bm{z}$若使$\Phi(\bm{z})$为真，就称它\term{满足}（satisfy）$\Phi$。
至少存在一个满足赋值时，公式\term{可满足}（satisfiable）；不存在时称为不可满足。
这里要判断的是
\[
  \exists\bm{z}\in\{0,1\}^n,\quad \Phi(\bm{z})=\mathrm{true},
\]
不是证明$\Phi$对每个赋值都真。对每个赋值都真的公式称为永真式；可满足性与永真性
对应不同的量词，不能混用。

\term{布尔可满足性问题}（Boolean satisfiability problem, SAT）的输入是一个布尔
公式，输出只回答它是否存在满足赋值。输入限制为CNF时称为CNF-SAT；进一步要求每个
子句至多含三个文字时称为\term{3-SAT}。若采用“恰含三个文字”的定义，可以重复文字
补齐较短子句，不改变可满足性。式~\eqref{eq:comb-running-cnf}因此既是一个CNF-SAT
实例，也是一个3-SAT实例。三者共享“是否存在满足赋值”这一判定接口，但允许的输入
语法不同，不能只凭名称互换实例。

逐项核对式~\eqref{eq:comb-running-cnf}可得五个满足集合，正好是
式~\eqref{eq:comb-running-feasible-family}中的$\mathcal{C}_0$。容斥从违反事件的
并集计算数量，逻辑公式则直接检验一个候选是否合格；两种表示回答的问题不同，但结果必须
一致。

\subsection{蕴含、多个模型与冲突}

若每个满足$\Phi$的赋值也满足命题$\Psi$，记作
\[
  \Phi\models\Psi,
\]
称$\Psi$是$\Phi$的\term{逻辑后承}（logical consequence）。这表示规则已经强制
推出$\Psi$，不表示某个求解器碰巧返回了满足$\Psi$的一个答案。例如，由
式~\eqref{eq:comb-running-cnf}可得
\[
  \Phi\models(z_3\Rightarrow\neg z_2),
\]
因为$z_3$推出$z_1$，而$z_1$与$z_2$互斥。还可得
\[
  \Phi\land z_4\land\neg z_2
  \models z_3\land z_1.
\]
后一结论使用了额外观测$z_4=1,z_2=0$，不能去掉这些条件后仍声称
$\Phi\models z_3$。事实上，$\Phi\land z_4\land\neg z_2$只有
$\{1,3,4\}$这一个模型；原公式$\Phi$则有五个模型。这组规则因此同时给出了唯一答案、
多个答案与下段冲突情形的可复核对照。

规则允许多个模型时，可满足性只确认“至少有一个答案”。$\{1\}$与$\{2,4\}$都满足
$\Phi$，所以规则本身没有在二者之间作出选择。若再加入$z_2$与$z_3$两个单位子句，
公式变为不可满足：$z_3$推出$z_1$，却又要求互斥的$z_2$为真。冲突表示全部硬规则
无法同时成立，不等于程序尚未找到答案。

判断$\Phi\models\Psi$可以转化为一次不可满足性检查：
\begin{equation}
  \Phi\models\Psi
  \quad\Longleftrightarrow\quad
  \Phi\land\neg\Psi\text{ 不可满足}.
  \label{eq:comb-entailment-unsat}
\end{equation}
若右侧可满足，其模型就是一个满足$\Phi$却违反$\Psi$的反例；若右侧不可满足，
这样的反例不存在。这个等价关系把全称式的逻辑后承转换为存在反例的搜索。

定义还包含一个容易误读的边界：若$\Phi$本身不可满足，则
$\operatorname{Mod}(\Phi)=\varnothing$，因而对任意$\Psi$都有
$\Phi\models\Psi$。这称为空真，并不表示冲突规则提供了关于$\Psi$的有效信息。
实际推理通常先检查$\Phi$可满足，再解释它的逻辑后承。

\subsection{从逻辑子句到零一约束}

布尔规则也可以写成关于$z_i\in\{0,1\}$的代数约束。一个文字若为$z_i$，
用数值$z_i$表示；若为$\neg z_i$，用$1-z_i$表示。子句
$\ell_1\lor\cdots\lor\ell_r$恰好等价于
\begin{equation}
  \sum_{j=1}^{r}v(\ell_j)\geq1,
  \qquad
  v(z_i)=z_i,\quad v(\neg z_i)=1-z_i.
  \label{eq:comb-clause-linearization}
\end{equation}
因为左侧计算为真的文字个数，至少一个文字为真当且仅当该和不小于$1$。

式~\eqref{eq:comb-running-cnf}因而等价于
\begin{equation}
  z_1+z_2=1,\qquad
  z_3\leq z_1,\qquad
  z_4\leq z_2+z_3,\qquad
  \bm{z}\in\{0,1\}^4.
  \label{eq:comb-running-integer-constraints}
\end{equation}
第一条等式同时包含“至少一个”和“至多一个”。这个代数表示把规则接到第
\ref{chap:optimization}章的可行域语言，但变量的整数性仍然关键。

若把$\bm{z}\in\{0,1\}^4$放宽为$\bm{z}\in[0,1]^4$，得到一个连续松弛。
分数点
\[
  \bm{z}=(1/2,1/2,1/2,1)
\]
满足全部松弛约束，却不是任何逻辑赋值。按阈值$1/2$把各坐标都取为$1$，还会违反
$z_1+z_2=1$。松弛点可提供优化界或帮助构造整数解，但它本身不是规则系统的合法答案；
取整后必须重新核对原约束。

\subsection{硬约束与软代价}

\term{硬约束}（hard constraint）直接定义可行集合，违反任意一条就不是候选答案。
有些规则表达偏好而非绝对禁止，此时可以为违反规则设置非负代价。令
$r_j(\bm{z})\in\{0,1\}$表示第$j$条规则是否被违反，权重$\lambda_j\geq0$，
则可以最小化
\begin{equation}
  c(\bm{z})+\sum_{j=1}^{m}\lambda_j r_j(\bm{z}),
  \qquad
  \bm{z}\in\{0,1\}^n.
  \label{eq:comb-soft-constraint-objective}
\end{equation}
$c$是任务本身的代价，第二项为规则冲突计分。权重有限时，足够大的任务收益可能抵消
违反规则的代价，所以软规则不再具有逻辑蕴含意义。

把$\lambda_j$取得很大也不能未经证明就等同于硬约束。若$c$的变化范围已知，例如
$|c(\bm{z})-c(\bm{z}')|\leq B$对任意赋值成立，并且存在一个不违反任何软规则的
赋值$\bm{z}_0$，那么令每个$\lambda_j>B$可以排除所有带违反项的全局最优解：
任一违规赋值的目标值严格大于$c(\bm{z}_0)$。若不存在零违反赋值、无法给出可用的
$B$，或希望一族不同规模的实例共用罚权而$c$的振幅没有统一上界，或者只把部分权重
取得较大，这个论证便不成立。
硬约束、罚函数和带松弛变量的约束各自改变了可行性或比较准则，使用时应明确允许返回
什么答案。

\section{组合优化与次模性}
\label{sec:comb-optimization-submodularity}

\subsection{可行性、最优性与求解代价}

规则确定集合族$\mathcal{C}\subseteq\mathcal{P}(V)$后，仍需用目标函数比较不同
可行集合。一个有限组合优化问题可写成
\begin{equation}
  \max_{S\in\mathcal{C}}F(S)
  \qquad\text{或}\qquad
  \min_{S\in\mathcal{C}}C(S).
  \label{eq:comb-discrete-optimization}
\end{equation}
由于$\mathcal{C}$有限且非空，实值目标一定达到最优值。这只回答存在性；验证
$S\in\mathcal{C}$、计算$F(S)$以及在不枚举全部候选的情况下找到最优解，仍是不同任务。
存在性的证明只需把$\mathcal{C}=\{S_1,\ldots,S_N\}$列成有限序列：有限个实数
$F(S_1),\ldots,F(S_N)$可逐项比较并取得最大值，达到该值的某个$S_i$就是最优解。
最小化结论可对$-C$应用同一论证。
若$\mathcal{C}=\varnothing$，则没有可行解；若目标允许取无穷值，也需另外说明所用的
大小关系。因此“有限”“非空”和“实值”都是这个存在性结论的条件。

继续使用四个候选。设需求集合为$R=\{a,b,c,d,e\}$，各候选覆盖
\[
  Q_1=\{a,b\},\quad
  Q_2=\{b,c,d\},\quad
  Q_3=\{c,e\},\quad
  Q_4=\{d,e\}.
\]
定义覆盖收益
\begin{equation}
  F(S)=\left|\bigcup_{i\in S}Q_i\right|.
  \label{eq:comb-running-coverage}
\end{equation}
在$\mathcal{C}_0$的五个可行集合上，收益依次为
$2,4,5,3,4$，所以$\{1,3,4\}$达到最优值$5$。这里最优解可以直接列举得到，
因为例子只有五个候选集合；若$n$增大，$2^n$个赋值的枚举时间很快就会成为主要障碍。

只检查局部交换也不能替代全局比较。设可行集合是
$V=\{1,2,3,4\}$的所有二元子集，并规定
$F(\{1,2\})=1$、$F(\{3,4\})=3$，其余二元子集的值均为$0$。
从$\{1,2\}$出发，替换其中任一元素都会使目标降为$0$，所以它是一交换局部最优解；
全局最优解却是$\{3,4\}$。局部最优能否推出全局最优，必须来自目标与可行域的额外结构。

近似解的误差也要先约定口径。对非负最大化目标，若算法输出$\widehat S$满足
$F(\widehat S)\geq\alpha F(S^\star)$，称其具有$\alpha$乘法近似，其中
$0\leq\alpha\leq1$；若满足
$F(S^\star)-F(\widehat S)\leq\varepsilon$，则是加法误差不超过$\varepsilon$。
当最优值为$0$或目标可取负值时，乘法比例可能失去区分力，加法误差或经过说明的平移、
归一化口径更合适。

若每个候选还有成本$c_i\geq0$，可以最大化
$F(S)-\sum_{i\in S}c_i$，也可以在预算
$\sum_{i\in S}c_i\leq B$下最大化$F(S)$。这两个问题通常不等价：固定罚权隐含一种
收益与成本交换率，预算约束则禁止所有超额集合。要断言对给定预算存在某个乘子并使
两类问题的最优解对应，需要第~\ref{chap:optimization}章所述的额外对偶条件。
一般离散问题不保证这样的乘子存在，不能只凭形式相似互换两类问题。

\subsection{边际收益递减}

一般集合函数可以任意指定每个子集的值，几乎没有可利用的结构。
\term{次模函数}（submodular function）要求已选择的集合越大，再加入同一元素所得的
边际收益越小。正式地，$F:\mathcal{P}(V)\to\mathbb{R}$若对所有
$A\subseteq B\subseteq V$和$e\in V\setminus B$满足
\begin{equation}
  F(A\cup\{e\})-F(A)
  \geq
  F(B\cup\{e\})-F(B),
  \label{eq:comb-diminishing-returns}
\end{equation}
就称$F$为次模函数。记
\[
  \Delta(e\mid A)=F(A\cup\{e\})-F(A)
\]
后，式~\eqref{eq:comb-diminishing-returns}表示
$\Delta(e\mid A)\geq\Delta(e\mid B)$。

次模性不等于\term{单调性}（monotonicity）。单调函数要求
$A\subseteq B$时$F(A)\leq F(B)$，也就是所有边际收益非负；次模函数允许边际收益
为负，只要求它随上下文扩大而不增加。模函数
$F(S)=c+\sum_{i\in S}w_i$的边际收益与上下文无关，因此同时是次模和超模函数，
不论某些$w_i$是否为负。

互补收益给出次模性失效的最小反例。令
$G(S)=\mathbb{I}\{\{1,3\}\subseteq S\}$。把元素$3$加入空集的边际收益为$0$，
把它加入$\{1\}$的边际收益却为$1$。由于$0<1$，边际收益随上下文增大而上升，
$G$不是次模函数。这个例子区分了“两个元素共同出现才有收益”的互补关系与覆盖函数中的
替代关系。

次模性还有一个对称的交并形式。

\begin{theorem}[次模性的两个等价定义]
\label{thm:comb-submodular-equivalence}
集合函数$F:\mathcal{P}(V)\to\mathbb{R}$满足边际收益递减
式~\eqref{eq:comb-diminishing-returns}，当且仅当对任意$A,B\subseteq V$有
\begin{equation}
  F(A)+F(B)\geq F(A\cup B)+F(A\cap B).
  \label{eq:comb-submodular-lattice}
\end{equation}
\end{theorem}

\begin{proof}
先假设边际收益递减。把$A\setminus B$中的元素依次记为$e_1,\ldots,e_r$。
从$A\cap B$加入这些元素得到$A$，从$B$加入同一序列得到$A\cup B$。
在第$j$步，前一条路径的当前集合包含于后一条路径的当前集合，因此
\[
  F(A)-F(A\cap B)
  \geq F(A\cup B)-F(B).
\]
移项即得式~\eqref{eq:comb-submodular-lattice}。

反过来，取$A\subseteq B$和$e\notin B$，将交并不等式用于
$A\cup\{e\}$与$B$。二者的交为$A$，并为$B\cup\{e\}$，故
\[
  F(A\cup\{e\})+F(B)
  \geq F(B\cup\{e\})+F(A).
\]
移项便得到式~\eqref{eq:comb-diminishing-returns}。
\end{proof}

\subsection{覆盖函数为何次模}

式~\eqref{eq:comb-running-coverage}是次模性的典型来源。更一般地，为每个需求
$r\in R$给定权重$w_r\geq0$，定义
\begin{equation}
  F(S)
  =
  \sum_{r\in\bigcup_{i\in S}Q_i}w_r.
  \label{eq:comb-weighted-coverage}
\end{equation}

\begin{theorem}[非负加权覆盖函数单调且次模]
\label{thm:comb-coverage-submodular}
若所有$w_r\geq0$，则式~\eqref{eq:comb-weighted-coverage}定义的$F$
是单调次模函数，并满足$F(\varnothing)=0$。
\end{theorem}

\begin{proof}
空集合不覆盖任何需求，所以$F(\varnothing)=0$。若$A\subseteq B$，则$A$覆盖的
需求也被$B$覆盖；权重非负保证$F(A)\leq F(B)$。对
$A\subseteq B$和$e\notin B$，加入$e$的边际收益是
\[
  \Delta(e\mid A)
  =
  \sum_{r\in Q_e\setminus\bigcup_{i\in A}Q_i}w_r.
\]
由于
$\bigcup_{i\in A}Q_i\subseteq\bigcup_{i\in B}Q_i$，在$B$的基础上尚未覆盖的
$Q_e$元素不会比在$A$的基础上更多。逐项使用$w_r\geq0$，得到
$\Delta(e\mid A)\geq\Delta(e\mid B)$。
\end{proof}

在贯穿例子中，候选$3$单独加入空集时新增$\{c,e\}$，边际收益为$2$；若已经选择
$\{2\}$，需求$c$已经覆盖，加入$3$只新增$e$，边际收益降为$1$。这个计算展示了
式~\eqref{eq:comb-diminishing-returns}的一种情形，但一般结论来自上面的集合包含
证明，而不是来自一个数值例子。

非负权重是单调性证明的实质条件。若某个被覆盖对象具有负权重，增加候选可能降低目标；
函数是否仍次模要按具体形式检查。两个次模函数之和仍次模，非负倍数也保持次模性，
但次模函数之差一般不次模。给覆盖收益减去模块化成本仍是次模函数，却可能失去单调性。

\subsection{基数约束下的贪心保证}

次模性约束边际收益之间的关系，但只有与特定可行域和算法结合，才能得到求解保证
\scite{math-nemhauser1978submodular}
考虑归一化单调次模最大化
\[
  \max_{S\subseteq V,\ |S|\leq k}F(S),
  \qquad F(\varnothing)=0.
\]
贪心算法从$S^{(0)}=\varnothing$开始，在第$t$步选择边际收益最大的元素
\[
  e^{(t)}\in\argmax_{e\in V\setminus S^{(t)}}\Delta(e\mid S^{(t)}),
  \qquad
  S^{(t+1)}=S^{(t)}\cup\{e^{(t)}\}.
\]
平局可按固定顺序处理，使算法输出确定。

\begin{theorem}[单调次模最大化的贪心界]
\label{thm:comb-greedy-submodular}
设$1\leq k\leq|V|$，$F$归一化、单调且次模，$S^\star$是在
$|S|\leq k$下的最优解。
上述贪心算法运行$k$步得到$S^{(k)}$，则
\begin{equation}
  F(S^{(k)})
  \geq
  \left(1-\left(1-\frac1k\right)^k\right)F(S^\star)
  \geq(1-e^{-1})F(S^\star).
  \label{eq:comb-greedy-bound}
\end{equation}
\end{theorem}

\begin{proof}
固定第$t$步。由单调性、次模性以及逐个加入$S^\star\setminus S^{(t)}$中的元素，
\begin{align*}
  F(S^\star)-F(S^{(t)})
  &\leq F(S^{(t)}\cup S^\star)-F(S^{(t)})\\
  &\leq\sum_{e\in S^\star\setminus S^{(t)}}\Delta(e\mid S^{(t)}).
\end{align*}
右侧至多有$k$项，所以其中至少一项不小于
$[F(S^\star)-F(S^{(t)})]/k$。若右侧没有元素，则单调性已经给出
$F(S^{(t)})\geq F(S^\star)$，所需递推平凡成立；否则贪心选择的边际收益不小于上述某一项，
故
\[
  F(S^{(t+1)})-F(S^{(t)})
  \geq\frac{F(S^\star)-F(S^{(t)})}{k}.
\]
令剩余差距$r^{(t)}=F(S^\star)-F(S^{(t)})$，便有
$r^{(t+1)}\leq(1-1/k)r^{(t)}$。由$F(S^{(0)})=0$递推得到
\[
  r^{(k)}\leq\left(1-\frac1k\right)^kF(S^\star),
\]
移项给出第一个不等式。再用
$(1-1/k)^k\leq e^{-1}$得到第二个不等式。
\end{proof}

定理~\ref{thm:comb-greedy-submodular}不能直接套到
式~\eqref{eq:comb-running-feasible-family}。逻辑规则形成的$\mathcal{C}_0$
不是普通基数约束，甚至不是向下封闭的：$\{1,3\}$可行，其子集$\{3\}$却不可行。
逐个加入当前边际收益最大的元素可能在中途违反规则，或走入无法扩展到最优解的分支。
其他约束下可以有不同的贪心、连续松弛或近似结果，但保证必须重新证明，不能只凭目标
具有次模性。

\subsection{连续松弛、界与离散答案}

把整数可行集$\mathcal{C}$嵌入$\{0,1\}^n$后，可以用一个较大的连续集合
$P\subseteq[0,1]^n$包住它。对最大化问题，若连续目标在整数点上与原目标一致，
松弛最优值通常给出原最优值的上界；任何取整后可行解则给出下界。对最小化问题方向相反。
这与第~\ref{sec:opt-relaxations-certificates}节的上下界口径一致。

例如考虑
\[
  \max\ z_1+z_2
  \quad\text{s.t.}\quad
  2z_1+2z_2\leq3,\qquad z_1,z_2\in\{0,1\}.
\]
整数最优值为$1$。放宽为$0\leq z_1,z_2\leq1$后，
$(z_1,z_2)=(1,1/2)$可取得松弛最优值$3/2$，确实是整数最优值的上界。
若把正坐标都取为$1$，则得到$(1,1)$并违反容量约束；若把$1/2$向下取整，则恢复一个
值为$1$的可行解。这个算例同时展示了松弛界、取整可行性与取整后目标值是三项不同检查。

即使松弛已经精确描述整数点的凸包，求得的分数最优点也可能位于一个最优面内部。
线性目标保证同一最优面上至少有整数极点，不表示任意返回的分数点都能逐坐标取整。
若松弛比凸包更大，还可能出现\term{整数间隙}（integrality gap）：松弛最优值严格优于
所有整数可行解。因而应分别记录松弛值、取整后的可行值及二者差距。

集合函数的连续扩展也有多种选择。无论采用何种扩展，连续变量上的凸性或凹性都不能仅由
“原函数次模”一句话替代；扩展方式决定可用的优化工具。第~\ref{chap:graph-theory}章会
讨论一个可精确求解的特殊情形：二值成对能量满足相应次模不等式时，可以表示为非负容量
的割。这里建立的一般集合函数定义本身并不保证任意次模最大化或最小化都能用图割解决。

\section{可计算性与复杂度}
\label{sec:comb-computability-complexity}

\subsection{先规定输入怎样交给程序}

复杂度不是对象本身附带的标签，而是相对于输入编码、计算模型和所求输出定义的。
对有限变量问题，可以把$n$、CNF公式、整数或有理数权重以及阈值写成有限比特串。
输入长度记为$L$，它包括变量索引、子句、数值的二进制位数和分隔信息。
若一个权重需要$b$位表示，读取和比较它就不能无条件算作与$b$无关的一步。

同一优化任务常对应三个不同计算问题。
\begin{itemize}
  \item \term{搜索问题}要求返回一个满足$\Phi$的$\bm{z}$；
  \item \term{优化问题}要求返回目标最优的可行$\bm{z}$或最优值；
  \item \term{判定问题}只问是否存在满足$\Phi$且$F(S(\bm{z}))\geq K$的赋值。
\end{itemize}
复杂性类别通常先为判定问题定义。若能多次调用判定算法并固定变量，常可恢复一个解；
但调用次数、数值精度和自归约结构都要计入，不能把“知道答案存在”直接当作已经构造答案。

\subsection{可枚举、可判定与可计算}

一个对象集合若存在程序依次输出其中全部对象，称为\term{可枚举的}
（computably enumerable）。这里要求程序只输出该集合中的对象，并最终输出其中每个
对象；程序可以重复输出，也不必告诉我们某个尚未出现的对象
以后是否会出现。一个判定问题若存在程序对每个合法输入都在有限步内停机，并正确回答
“是”或“否”，称为\term{可判定的}（decidable）。若程序只保证在“是”实例上最终
停机并接受，而在“否”实例上可能永不停止，则称该问题\term{半可判定的}
（semidecidable）。
类似地，字符串上的函数$f$若存在一个程序，对定义域内每个输入$x$都停机并输出
$f(x)$，就称$f$为\term{可计算函数}（computable function）。判定问题可看作计算
取值为$0$或$1$的特征函数，因此可判定性比只确认正例的半可判定性要求更强。

可枚举不等于可判定。若某集合及其补集都可枚举，可以交错运行两个枚举或半判定程序：
输入最终会被其中一方确认，于是得到判定程序。若只有正例可枚举，等待尚未出现的对象
无法证明它永远不会出现。

在明确的有限编码下，CNF可满足性是可判定的。按字典序枚举全部
$2^n$个赋值，每个赋值逐子句检查；找到满足赋值就回答“是”，全部检查完仍未找到则
回答“否”。若公式有$m$个子句、总文字出现次数为$q$，直接实现需要
$O(2^n q)$量级的时间和存放当前赋值及公式所需的空间。这个算法很慢，但一定停机，
所以它证明了可判定性。

“每个实例的候选集有限”还不足以独立构成上述证明。程序必须能从输入有效得到候选编码，
能判定一个候选是否满足约束，并能知道枚举何时完成。若把一个不可计算实数或一个可能
不停机的黑盒藏进“目标求值”，有限个候选也未必能被统一程序比较。可计算性要求的是
对整族输入存在同一个有限描述的程序。

\subsection{存在不可判定的问题}

有限变量问题在有效表示下可以穷举，但一般程序的行为并非总能判定。
令
\[
  \mathrm{HALT}
  =
  \{(P,x):\text{程序 }P\text{ 在输入 }x\text{ 上最终停机}\}.
\]
它是半可判定的：模拟$P(x)$，若模拟停止就接受；若$P(x)$永不停止，模拟也不会给出
否定答案。下面的对角论证说明不存在总能判断停机与否的程序。

停机问题的不可判定性是可计算性理论的基本结论\scite{math-turing1937computable}
\begin{theorem}[停机问题不可判定]
\label{thm:comb-halting-undecidable}
$\mathrm{HALT}$不可判定。
\end{theorem}

\begin{proof}
反设存在判定程序$H(P,x)$，它总会停机，并在$P(x)$停机时返回$1$，否则返回$0$。
据此构造程序$D$：输入一个程序编码$P$，先计算$H(P,P)$；若结果为$1$，就进入无限
循环；若结果为$0$，就立即停机。

现在把$D$自己的编码交给$D$。若$H(D,D)=1$，按照$D$的定义，$D(D)$进入无限循环，
与$H$判断它会停机矛盾。若$H(D,D)=0$，$D(D)$立即停机，又与$H$判断它不停机矛盾。
两种返回值都导致矛盾，所以这样的判定程序$H$不存在。
\end{proof}

定理~\ref{thm:comb-halting-undecidable}区分了“某次模拟尚未结束”和“已经证明不会结束”。
前者是有限时刻的观察，后者要求排除所有未来时刻。不可判定也不等于每个具体程序都无法
分析；许多程序可由终止度量、循环不变量或有限状态搜索证明停机。结论是不存在一个对
任意程序和输入都正确的统一判定器。

\subsection{多项式时间、证书与复杂性类别}

可判定问题仍可能需要随输入长度急剧增长的资源。若一个确定性算法的最坏运行时间由
$L$的某个固定次数多项式控制，该判定问题属于复杂性类别$\mathrm{P}$。
类别$\mathrm{NP}$包含满足如下条件的判定问题$\mathcal{L}$：存在多项式$p$和确定性
多项式时间验证程序$\mathsf{Ver}$，使每个输入$x$都满足
\[
  x\in\mathcal{L}
  \quad\Longleftrightarrow\quad
  \exists y,\quad |y|\leq p(|x|),\quad \mathsf{Ver}(x,y)=1.
\]
其中$|x|$是输入的比特长度，$y$称为\term{证书}（certificate）。双向条件不仅要求
每个“是”实例存在可接受的多项式长度证书，还要求“否”实例的所有这类证书都被拒绝。
名称中的“非确定性”描述这种猜测证书后验证的计算模型，不表示算法可以容忍随机错误。
确定性多项式算法的运行轨迹本身可作证书，所以
$\mathrm{P}\subseteq\mathrm{NP}$。目前并不知道这个包含是否严格；把
$\mathrm{P}\neq\mathrm{NP}$当作已经证明的事实，会超出这里可用的结论。

CNF可满足性属于$\mathrm{NP}$。证书就是$n$位赋值$\bm{z}$，验证程序逐个检查子句，
时间关于公式长度为线性量级。穷举证明它可判定，短证书证明它属于$\mathrm{NP}$；
二者都没有给出多项式时间的求解算法。

若问题$A$的每个实例$x$都能在多项式时间内变换为问题$B$的实例$f(x)$，并满足
\[
  x\in A\quad\Longleftrightarrow\quad f(x)\in B,
\]
就称$A$可\term{多项式时间多对一归约}（polynomial-time many-one reduction）
到$B$，记作$A\leq_{\mathrm{p}}B$。归约方向表示：若能高效解决$B$，就能先计算
$f(x)$再高效解决$A$。因此从已知困难问题归约到新问题，才能证明新问题至少同样困难；
把方向写反不会得到这个结论。

一个问题若所有$\mathrm{NP}$问题都可多项式时间归约到它，就称为
\term{NP困难}（NP-hard）；若它还属于$\mathrm{NP}$，就称为
\term{NP完全}（NP-complete）。Cook--Levin定理给出SAT是NP完全问题
\scite{math-cook1971complexity}通过引入
辅助变量，可以把一般公式在多项式时间内变成等可满足的CNF公式；再把长子句拆成由辅助
变量连接的三文字子句，可以归约到3-SAT。这里的“等可满足”只保证原公式有解当且仅当
新公式有解，不要求辅助变量加入后两个公式拥有相同的模型集合。由这些接口可知CNF-SAT
与3-SAT也都是NP完全问题。这里把结论作为复杂度理论的基准，不展开对任意多项式时间
验证计算的编码证明；下面只证明与本章约束表示直接相关的归约。

\begin{theorem}[CNF可满足性归约到零一线性可行性]
\label{thm:comb-sat-to-binary-linear}
给定任意CNF公式，可以在关于公式长度的多项式时间内构造一组零一线性不等式，
使该公式可满足当且仅当这些不等式存在零一可行解。因此零一线性可行性是NP困难的；
当所有整数系数按二进制编码时，它也属于NP，因而是NP完全的。
\end{theorem}

\begin{proof}
为每个命题变量建立$z_i\in\{0,1\}$。对每个子句
$C=\ell_1\lor\cdots\lor\ell_r$，加入式~\eqref{eq:comb-clause-linearization}：
\[
  \sum_{j=1}^{r}v(\ell_j)\geq1.
\]
每个文字只产生一个系数为$1$的变量项或$1-z_i$，所以构造时间和输出长度都关于原公式
长度为线性量级。

若原公式有满足赋值，把真假分别编码为$1,0$。每个子句至少一个文字为真，对应不等式
左侧至少为$1$，所以得到零一可行解。反过来，任一零一可行解定义一个布尔赋值；
每条不等式左侧至少为$1$，说明相应子句至少一个文字为真，故全部子句同时满足。
两边可行性等价，归约成立。

由CNF可满足性的NP困难性，零一线性可行性NP困难。对其NP成员关系，证书是一组零一
变量值；验证者按二进制整数运算检查所有不等式，所需时间关于输入和证书长度为多项式。
\end{proof}

定理~\ref{thm:comb-sat-to-binary-linear}说明，换成代数记号不会自动消除逻辑搜索的
困难。它也没有声称每个零一线性实例都难：式~\eqref{eq:comb-running-integer-constraints}
只有四个变量，可以立即求解；某些约束矩阵或关系结构还允许专门的多项式算法。
NP完全描述的是一族问题在最坏情况下的归约地位，不是每个实例的运行时间预测。

\subsection{复杂度结论的边界}

输入规模必须与算法成本使用同一口径。枚举$2^n$个赋值是关于变量数$n$的指数时间；
若输入本身已经显式列出全部$2^n$个候选，那么扫描列表反而关于输入长度是线性的。
把隐式表示与显式列表混在一起，会让同一算法同时被误称为指数和线性。

数值大小也不能替代编码长度。若整数$N$用二进制给出，输入它只需
$\lfloor\log_2N\rfloor+1$位；执行$N$轮的循环关于数值$N$是线性的，却关于输入位数
是指数级的。运行时间关于数值大小是多项式、关于二进制输入长度却未必是多项式的算法，
称为\term{伪多项式时间算法}（pseudo-polynomial-time algorithm）。只有同时说明
数值编码与所计基本运算，才能判断一个界是否真正关于输入长度为多项式。

最坏情况复杂度也不同于一次实验运行时间。某个求解器在当前数据上很快，可能来自实例
结构、预处理或启发式分支顺序；这不能证明问题属于$\mathrm{P}$。反过来，NP完全性
不排除小实例、特殊结构、参数受限情形或近似目标能够高效求解。应说明保证针对精确解、
近似值还是可行解，并列明依赖的参数与表示。

统计样本量与计算时间同样不能互相替代。一个候选族可能只需较少样本就能控制泛化误差，
经验风险最小化却仍难以计算；也可能存在高效优化算法，但样本不足以区分多个预测规则。
可计算学习还要求训练程序和输出的预测程序都停机，高效学习进一步要求相应资源由规定
参数的多项式控制。组合计数可以进入统计界，也可以进入枚举成本，但两处的规模对象和
结论必须分别核对。

\section{本章小结}
\label{sec:comb-summary}

有限离散问题至少包含四层约束。集合族与计数说明候选空间有多大；容斥校正重叠事件，
Möbius反演则从所有子集的累计量恢复局部贡献。二者共享交替抵消机制，但反演本身只提供
代数重构，不会自动给出概率因子分解或非负局部项。

逻辑公式规定哪些赋值允许。可满足表示至少存在一个模型，逻辑后承要求所有模型都满足
结论，不可满足则表示硬规则冲突。把子句改写为零一不等式可以接入优化工具，连续松弛和
软代价却改变了答案对象；分数解、取整结果和硬约束模型必须分开核对。

目标函数在可行集合之间作选择。次模性刻画边际收益递减，单调性另行要求增加元素不降低
目标。归一化单调次模函数在基数约束下具有贪心近似界，但一般逻辑约束、非单调目标或其他
松弛都需要新的算法条件。候选有限保证最优值存在，不保证无需枚举就能找到最优解。

计算层面还要依次区分可判定、可计算和多项式时间。穷举证明有效表示下的有限CNF问题
可判定，停机问题说明一般程序性质可能不可判定，归约则比较可判定问题的最坏情况难度。
经验运行时间不能替代复杂度证明，NP完全性也不排除特殊结构上的高效算法。第
\ref{chap:graph-theory}章将把本章变量之间的关系表示为图，研究分解、路径与局部结构
怎样改变计算组织。
